% Programming with InitechBase, Chapter 4: Procedures and Parameters
% (AN-16; annex Sec 8.2 P4).  Panels replayed (examples/p4-*.dot); sources
% cmd/proc.c (PROCEDURE, DO ... WITH by value, PARAMETERS, RETURN, PUBLIC),
% samir_main.c (DO of a disk program only at the dot prompt).
\chapter{Procedures and Parameters}
\Gates{test-pubs-examples, test-interp-proc}
A calculation that a program needs more than once, or that makes the program
hard to read, is written once as a procedure and called by name.  This chapter
shows you how to write a procedure, pass values to it and get its result back.

\begin{ChapterCovers}
\item Procedures in a program file
\item Calling a procedure with \cmd{DO}
\item Passing values with \kw{WITH} and \cmd{PARAMETERS}
\item Returning a result
\end{ChapterCovers}

\ATEHead{Preparing for This Chapter}
Start InitechBase with a fresh copy of the sample diskette.

\ATEHead{Procedures}
The procedures of a program follow its main program in the same file.  Each
begins with a \cmd{PROCEDURE} line that names it and ends with \cmd{RETURN}.
The main program ends at the first \cmd{PROCEDURE} line, so a procedure runs only
when \cmd{DO} calls it.

A program calls the procedures of its own file.  A program file is run from the
dot prompt; one program does not run another program file.  The main menu of
the system (Chapter~11) therefore holds its procedures in its own file.

\ATEHead{Passing Values}
\kw{WITH} after the name of the procedure gives the values to pass, separated by
commas.  The first command of the procedure, \cmd{PARAMETERS}, names the memory
variables that receive them, in the same order.  The procedure receives copies:
what it does with its parameters does not change the values in the caller.

\file{LATEFEE.PRG} computes the late fee of an invoice: one per cent of the
amount for each full thirty days the invoice is past due.  The procedure
\texttt{FEE} receives the amount and the number of days:
\ProgramListing{LATEFEE.PRG}

\ATEHead{Returning a Result}
A procedure gives back its result in a memory variable declared \cmd{PUBLIC},
which remains when the procedure returns.  \texttt{FEE} stores the fee in
\texttt{MFEE}:
\begin{Steps}
\item Type \type{USE INVOICE}, then \type{GO 2}.
\item Type \type{DO LATEFEE}, then \type{? MFEE}.
\end{Steps}
\DotExample{examples/p4-latefee}
The parameters \texttt{PAMOUNT} and \texttt{PDAYS} are private to the procedure
and are released when it returns.

\begin{AteNote}
Pass character values to a procedure through memory variables or expressions
(\doctitle{Using InitechBase}, Chapter~5, \cmd{DO}).
\end{AteNote}

\ATEHead{Where Do You Go From Here?}
Chapter~5 looks at the functions that build the values programs print.

\ATEHead{Summary}
\begin{enumerate}
\item A procedure begins with \cmd{PROCEDURE} \xp{name} and ends with
  \cmd{RETURN}; it follows the main program in the file.
\item \cmd{DO} \xp{procedure} \kw{WITH} \xp{values} calls it;
  \cmd{PARAMETERS} \xp{names} receives the values.
\item A result is returned in a \cmd{PUBLIC} memory variable.
\end{enumerate}
